\documentclass[../DoAn.tex]{subfiles}
\begin{document}

\label{chap:danh_gia_thuc_nghiem}

Chương này trình bày thiết kế thực nghiệm và phân tích kết quả của hệ thống phát hiện drone dựa trên tín hiệu RF. Các thí nghiệm được xây dựng nhằm đánh giá hai chiến lược sử dụng mô hình tiền huấn luyện là \textit{linear probing} và \textit{fine-tuning}, đồng thời khảo sát khả năng chuyển giao của năm backbone thị giác gồm ResNet50, EfficientNet-B2, ConvNeXtV2, SigLIP và DINOv2 sang miền ảnh spectrogram. Ngoài phép đánh giá trên bộ dữ liệu DroneDetect v2, các mô hình còn được kiểm tra trên hai tập dữ liệu RF tự thu bằng bladeRF tại băng tần 2,4~GHz. Cách tổ chức này cho phép xem xét đồng thời hiệu năng trong miền dữ liệu chuẩn, mức độ tổng quát hóa sang điều kiện thu mới và ảnh hưởng của chất lượng spectrogram đến kết quả nhận dạng.

Các nhận xét trong chương được giới hạn theo số liệu thực nghiệm hiện có. Những bảng liên quan đến cấu hình phần cứng, tham số huấn luyện và kết quả trên DroneDetect v2 được tạo dưới dạng khung LaTeX để bổ sung số liệu chính xác sau khi tổng hợp nhật ký thí nghiệm. Đối với dữ liệu tự thu, các tỷ lệ nhận dạng đã xác định được trình bày trực tiếp và được quy đổi sang số mẫu tương ứng trên mỗi tập gồm 200 spectrogram. Chương không sử dụng kết quả của một kiến trúc riêng lẻ để suy rộng thành kết luận cho toàn bộ nhóm CNN hoặc Transformer, bởi năm backbone khác nhau đồng thời về kiến trúc, quy mô và phương pháp tiền huấn luyện.

\section{Mục tiêu và thiết kế thực nghiệm}

Thực nghiệm được tổ chức nhằm so sánh hai chiến lược khai thác mô hình tiền huấn luyện và đánh giá mức độ phù hợp của từng backbone đối với ảnh spectrogram RF. Ở chiến lược linear probing, backbone được đóng băng và SVM được huấn luyện trên embedding trích xuất từ ảnh. Ở chiến lược fine-tuning, lớp phân loại cuối được thay thế và toàn bộ tham số backbone được cập nhật theo dữ liệu DroneDetect v2. Việc áp dụng hai chiến lược trên cùng các backbone giúp xác định lợi ích của quá trình thích nghi đặc trưng với miền RF.

Năm backbone được khảo sát gồm ResNet50, EfficientNet-B2, ConvNeXtV2, SigLIP và DINOv2. Hiệu năng trên DroneDetect v2 được đánh giá bằng accuracy, precision, recall, F1-score và ma trận nhầm lẫn. Các chỉ số này cho phép xem xét đồng thời tỷ lệ dự đoán đúng, khả năng hạn chế cảnh báo sai và khả năng tránh bỏ sót drone. Phép so sánh được thực hiện riêng cho từng backbone để phân biệt ảnh hưởng của chiến lược huấn luyện với ảnh hưởng của kiến trúc.

Bên cạnh dữ liệu benchmark, các mô hình được suy luận trên Drone1 và Drone2, là hai tập RF tự thu bằng bladeRF. Dữ liệu thực tế không được sử dụng để cập nhật tham số hoặc lựa chọn checkpoint. Thiết kế này nhằm kiểm tra khả năng duy trì hiệu năng khi thiết bị thu, mức nhiễu và điều kiện môi trường khác với dữ liệu huấn luyện. Kết quả còn được đối chiếu với chất lượng spectrogram để phân tích ảnh hưởng của SNR, clipping và stripe artifact.

Với backbone $m$, mức cải thiện accuracy của fine-tuning so với linear probing được tính bởi
\begin{equation}
    \Delta \mathrm{Acc}_{m}
    =
    \mathrm{Acc}^{\mathrm{FT}}_{m}
    -
    \mathrm{Acc}^{\mathrm{LP}}_{m},
    \label{eq:accuracy_gain}
\end{equation}
trong đó $\mathrm{Acc}^{\mathrm{FT}}_{m}$ và $\mathrm{Acc}^{\mathrm{LP}}_{m}$ lần lượt là accuracy của hai chiến lược. Giá trị dương cho thấy fine-tuning cải thiện chỉ số đang xét. Chênh lệch này được phân tích cùng precision, recall và F1-score thay vì sử dụng độc lập.

Để khảo sát thay đổi giữa benchmark và dữ liệu tự thu, recall của lớp drone trên DroneDetect v2 được đặt cạnh tỷ lệ nhận diện trung bình trên Drone1 và Drone2. Hai đại lượng có cùng ý nghĩa là tỷ lệ mẫu drone được phát hiện đúng. Việc sử dụng recall tránh so sánh trực tiếp accuracy của một tập hai lớp với tỷ lệ nhận diện của các tập chỉ chứa lớp drone.

\section{Dữ liệu và môi trường thực nghiệm}

\subsection{Bộ dữ liệu DroneDetect v2}

DroneDetect v2 là nguồn dữ liệu chính dùng để xây dựng và đánh giá mô hình trong miền benchmark. Dữ liệu ban đầu chứa tín hiệu RF được lưu dưới dạng I/Q cùng thông tin nhãn. Theo phương pháp ở Chương~3, các chuỗi I/Q được kiểm tra, phân đoạn và chuyển thành spectrogram bằng STFT. Mỗi spectrogram kế thừa nhãn và mã nguồn của bản ghi ban đầu để hỗ trợ chia tập và truy vết dữ liệu.

Trong phạm vi đề tài, các nhãn gốc được quy về hai lớp. Lớp \textit{Drone} bao gồm các đoạn có hoạt động RF của drone, còn lớp \textit{Non-drone} bao gồm các đoạn không chứa hoạt động mục tiêu. Cách quy đổi phù hợp với mục tiêu phát hiện sự hiện diện của drone, đồng thời tránh yêu cầu mô hình phân biệt từng loại thiết bị khi số lượng và phân bố mẫu giữa các loại có thể không đồng đều.

Dữ liệu được chia thành tập huấn luyện, tập xác thực và tập kiểm thử. Tập huấn luyện dùng để tối ưu SVM hoặc tham số mạng trong fine-tuning. Tập xác thực dùng để lựa chọn siêu tham số và checkpoint. Tập kiểm thử chỉ được sử dụng sau khi quá trình lựa chọn mô hình hoàn tất. Các spectrogram sinh từ cùng một bản ghi hoặc phiên thu cần được giữ trong cùng một tập để hạn chế rò rỉ đặc trưng nguồn.

\begin{table}[H]
    \centering
    \caption{Thống kê các tập dữ liệu được xây dựng từ DroneDetect v2}
    \label{tab:dronedetect_statistics}
    \renewcommand{\arraystretch}{1.3}
    \begin{tabular}{|p{0.23\textwidth}|c|c|c|p{0.25\textwidth}|}
        \hline
        \textbf{Tập dữ liệu} & \textbf{Drone} & \textbf{Non-drone} & \textbf{Tổng số} & \textbf{Mục đích} \\ \hline
        Huấn luyện & -- & -- & -- & Tối ưu tham số mô hình \\ \hline
        Xác thực & -- & -- & -- & Chọn siêu tham số và checkpoint \\ \hline
        Kiểm thử & -- & -- & -- & Báo cáo kết quả benchmark \\ \hline
    \end{tabular}
\end{table}

Bảng~\ref{tab:dronedetect_statistics} cần được điền bằng số lượng ảnh thực tế sau bước loại bỏ mẫu lỗi và chia theo nguồn. Tỷ lệ giữa hai lớp có ảnh hưởng đến cách diễn giải accuracy. Nếu tập kiểm thử cân bằng, mức $50\%$ tương ứng với dự đoán một lớp cố định; nếu tập không cân bằng, cần đối chiếu thêm phân bố lớp, balanced accuracy hoặc macro F1-score trước khi đánh giá mô hình có tốt hơn mức nền hay không.

\subsection{Dữ liệu RF tự thu}

Dữ liệu thực tế được thu bằng thiết bị bladeRF trong băng tần 2,4~GHz và lưu dưới dạng chuỗi I/Q. Quy trình thu sử dụng ăng-ten phù hợp với băng tần quan sát và phần mềm điều khiển SDR. Các bản ghi sau đó được kiểm tra chất lượng, phân đoạn và chuyển thành spectrogram bằng pipeline tương ứng với dữ liệu benchmark. Dữ liệu tự thu không tham gia quá trình huấn luyện hoặc lựa chọn checkpoint.

Hai tập đánh giá được ký hiệu là Drone1 và Drone2. Mỗi tập gồm 200 ảnh spectrogram có chứa tín hiệu drone, tương ứng tổng cộng 400 mẫu. Việc tách thành hai tập cho phép quan sát mức ổn định của mô hình giữa hai lần thu hoặc hai điều kiện tín hiệu. Thông tin chi tiết về khoảng cách, độ lợi, tốc độ lấy mẫu và điều kiện môi trường cần được bổ sung từ nhật ký thu để hỗ trợ giải thích chênh lệch kết quả.

\begin{table}[H]
    \centering
    \caption{Thống kê dữ liệu RF tự thu}
    \label{tab:real_data_statistics}
    \renewcommand{\arraystretch}{1.3}
    \begin{tabular}{|p{0.19\textwidth}|c|c|p{0.39\textwidth}|}
        \hline
        \textbf{Tập} & \textbf{Số ảnh} & \textbf{Nhãn} & \textbf{Điều kiện thu} \\ \hline
        Drone1 & 200 & Drone & -- \\ \hline
        Drone2 & 200 & Drone & -- \\ \hline
    \end{tabular}
\end{table}

Do Drone1 và Drone2 chỉ chứa lớp \textit{Drone}, chỉ số thu được trên hai tập là tỷ lệ nhận diện đúng lớp drone, tương đương recall của lớp dương trong từng tập. Các tập này không cho phép tính precision, specificity hoặc tỷ lệ cảnh báo giả vì không có mẫu \textit{Non-drone}. Hạn chế này cần được giữ trong cách diễn giải: kết quả cao chứng minh khả năng nhận diện drone trong các phiên thu đã xét, nhưng chưa đủ để khẳng định hệ thống hạn chế tốt cảnh báo sai trong môi trường thực tế.

\subsection{Môi trường huấn luyện}

Môi trường thực nghiệm ảnh hưởng trực tiếp đến khả năng tái lập thời gian huấn luyện, kích thước batch và một số sai khác số học. Các thí nghiệm được triển khai bằng framework học sâu và thư viện học máy để thực hiện tiền xử lý, fine-tuning và huấn luyện SVM. Thông tin phần cứng và phiên bản phần mềm cần được ghi chính xác theo máy đã sử dụng.

\begin{table}[H]
    \centering
    \caption{Cấu hình môi trường huấn luyện và đánh giá}
    \label{tab:training_environment}
    \renewcommand{\arraystretch}{1.3}
    \begin{tabular}{|p{0.34\textwidth}|p{0.51\textwidth}|}
        \hline
        \textbf{Thành phần} & \textbf{Cấu hình} \\ \hline
        Hệ điều hành & -- \\ \hline
        CPU & -- \\ \hline
        GPU & -- \\ \hline
        RAM & -- \\ \hline
        Ngôn ngữ lập trình & Python -- \\ \hline
        Framework học sâu & PyTorch -- \\ \hline
        Thư viện mô hình & -- \\ \hline
        Thư viện học máy & scikit-learn -- \\ \hline
        Thư viện xử lý tín hiệu & NumPy, SciPy -- \\ \hline
        Thiết bị SDR & bladeRF -- \\ \hline
    \end{tabular}
\end{table}

Cấu hình huấn luyện cũng cần được báo cáo để phân biệt ảnh hưởng của backbone với ảnh hưởng của siêu tham số. Bảng~\ref{tab:training_configuration} tạo khung cho các tham số chung. Nếu từng backbone sử dụng cấu hình riêng, bảng có thể được mở rộng thành một cột cho mỗi mô hình.

\begin{table}[H]
    \centering
    \caption{Cấu hình huấn luyện và tiền xử lý}
    \label{tab:training_configuration}
    \renewcommand{\arraystretch}{1.3}
    \begin{tabular}{|p{0.37\textwidth}|p{0.45\textwidth}|}
        \hline
        \textbf{Tham số} & \textbf{Giá trị} \\ \hline
        Kích thước ảnh đầu vào & -- \\ \hline
        Kích thước FFT & -- \\ \hline
        Cửa sổ STFT & Hann \\ \hline
        Hop length & -- \\ \hline
        Batch size & -- \\ \hline
        Số epoch tối đa & -- \\ \hline
        Hàm mất mát & -- \\ \hline
        Bộ tối ưu & -- \\ \hline
        Learning rate & -- \\ \hline
        Learning-rate scheduler & -- \\ \hline
        Tiêu chí chọn checkpoint & -- \\ \hline
        Cấu hình SVM & -- \\ \hline
    \end{tabular}
\end{table}

\section{Các phương pháp so sánh và tiêu chí đánh giá}

\subsection{Các backbone được khảo sát}

ResNet50 là backbone CNN sử dụng các khối residual và kết nối tắt. Mô hình được chọn làm kiến trúc tham chiếu phổ biến cho bài toán thị giác. Trong spectrogram, các phép tích chập có thể nhận biết cấu trúc cục bộ như biên dải phổ, vùng năng lượng tập trung và burst ngắn theo thời gian.

EfficientNet-B2 sử dụng chiến lược compound scaling để cân bằng độ sâu, độ rộng và độ phân giải đầu vào. Mô hình có hiệu quả tham số tương đối tốt và cung cấp một điểm so sánh với ResNet50. Kết quả của EfficientNet-B2 còn giúp kiểm tra liệu kiến trúc CNN được tối ưu về quy mô có duy trì được đặc trưng hữu ích khi miền dữ liệu chuyển từ ảnh tự nhiên sang spectrogram hay không.

ConvNeXtV2 là CNN hiện đại với các khối tích chập được thiết kế theo nhiều nguyên tắc của kiến trúc thị giác thế hệ mới. Kernel lớn và cách tổ chức tầng cho phép mô hình tổng hợp ngữ cảnh rộng hơn so với CNN truyền thống. Việc đưa ConvNeXtV2 vào thực nghiệm giúp so sánh một CNN hiện đại với ResNet50 và EfficientNet-B2 trong cùng quy trình.

SigLIP sử dụng encoder ảnh được tiền huấn luyện trong mô hình thị giác--ngôn ngữ với hàm mất mát sigmoid. Backbone này đại diện cho nhóm đặc trưng được học từ quan hệ giữa ảnh và mô tả văn bản trên dữ liệu quy mô lớn. Thí nghiệm kiểm tra liệu biểu diễn đó có thể thích nghi với ảnh spectrogram, vốn không chứa ngữ nghĩa màu sắc và vật thể giống ảnh tự nhiên.

DINOv2 là backbone thị giác được tiền huấn luyện tự giám sát. Mô hình hướng tới việc học biểu diễn tổng quát mà không phụ thuộc trực tiếp vào nhãn lớp của bài toán đích. DINOv2 được khảo sát vì các đặc trưng tự giám sát có thể mô tả cấu trúc hình học và quan hệ giữa các vùng ảnh, từ đó có khả năng chuyển giao sang mẫu thời gian--tần số.

Mỗi backbone được đánh giá bằng cùng hai chiến lược linear probing và fine-tuning, với điều kiện dữ liệu và tiêu chí lựa chọn mô hình tương ứng. Việc so sánh được tiến hành ở mức từng mô hình. Không sử dụng kết quả của DINOv2 hoặc SigLIP để đại diện tuyệt đối cho mọi Transformer, cũng không sử dụng kết quả của ResNet50 để đại diện cho mọi CNN.

\subsection{Accuracy}

Accuracy là tỷ lệ số mẫu được dự đoán đúng trên tổng số mẫu:
\begin{equation}
    \mathrm{Accuracy}
    =
    \frac{TP+TN}{TP+TN+FP+FN},
    \label{eq:accuracy_ch4}
\end{equation}
trong đó $TP$ là số mẫu drone được dự đoán đúng, $TN$ là số mẫu non-drone được dự đoán đúng, $FP$ là số mẫu non-drone bị dự đoán thành drone và $FN$ là số mẫu drone bị bỏ sót.

Accuracy cung cấp đánh giá tổng thể khi phân bố hai lớp tương đối cân bằng. Nếu tập kiểm thử mất cân bằng, một mô hình dự đoán thiên về lớp chiếm đa số vẫn có thể đạt accuracy cao. Vì vậy, chỉ số này được phân tích cùng precision, recall, F1-score và ma trận nhầm lẫn.

\subsection{Precision}

Precision của lớp \textit{Drone} đo tỷ lệ mẫu thực sự là drone trong số các mẫu được mô hình dự đoán là drone:
\begin{equation}
    \mathrm{Precision}
    =
    \frac{TP}{TP+FP}.
    \label{eq:precision_ch4}
\end{equation}

Precision thấp cho thấy hệ thống tạo nhiều cảnh báo sai. Trong bối cảnh triển khai, cảnh báo sai liên tục có thể làm giảm độ tin cậy của hệ thống. Chỉ số này không thể tính trên Drone1 và Drone2 nếu hai tập không chứa mẫu \textit{Non-drone}, bởi khi đó không xác định được số lượng false positive.

\subsection{Recall}

Recall đo tỷ lệ drone thực tế được phát hiện đúng:
\begin{equation}
    \mathrm{Recall}
    =
    \frac{TP}{TP+FN}.
    \label{eq:recall_ch4}
\end{equation}

Recall phản ánh trực tiếp khả năng hạn chế bỏ sót drone. Với Drone1 và Drone2, toàn bộ 200 mẫu của mỗi tập mang nhãn \textit{Drone}; do đó tỷ lệ nhận diện được báo cáo chính là recall. Chẳng hạn, tỷ lệ $91\%$ trên 200 mẫu tương ứng 182 mẫu được phát hiện và 18 mẫu bị bỏ sót.

\subsection{F1-score}

F1-score là trung bình điều hòa của precision và recall:
\begin{equation}
    \mathrm{F1}
    =
    2
    \frac{
        \mathrm{Precision}\times\mathrm{Recall}
    }{
        \mathrm{Precision}+\mathrm{Recall}
    }.
    \label{eq:f1_ch4}
\end{equation}

Chỉ số này giảm mạnh khi một trong hai thành phần precision hoặc recall thấp, vì vậy phù hợp để đánh giá sự cân bằng giữa cảnh báo sai và bỏ sót. Đối với bài toán hai lớp, macro F1-score có thể được tính bằng trung bình F1-score của hai lớp để tránh lớp có nhiều mẫu chi phối kết quả. Bảng benchmark cần ghi rõ sử dụng F1-score của lớp drone hay macro F1-score.

\subsection{Confusion Matrix}

Ma trận nhầm lẫn mô tả phân bố dự đoán theo nhãn thật và cho biết trực tiếp số lượng TP, TN, FP và FN. So với một chỉ số tổng hợp, ma trận này giúp nhận biết mô hình thiên về lớp nào. Một mô hình có accuracy thấp do chủ yếu bỏ sót drone cần được xử lý khác với mô hình có cùng accuracy nhưng tạo nhiều cảnh báo sai.

\begin{table}[H]
    \centering
    \caption{Cấu trúc ma trận nhầm lẫn của bài toán phát hiện drone}
    \label{tab:confusion_matrix}
    \renewcommand{\arraystretch}{1.3}
    \begin{tabular}{|c|c|c|}
        \hline
        \textbf{Nhãn thực tế} & \textbf{Dự đoán Drone} & \textbf{Dự đoán Non-drone} \\ \hline
        Drone & TP & FN \\ \hline
        Non-drone & FP & TN \\ \hline
    \end{tabular}
\end{table}

Trong phần kết quả DroneDetect v2, ma trận nhầm lẫn của từng cấu hình tốt nhất cần được bổ sung dưới dạng hình hoặc bảng. Việc báo cáo số lượng tuyệt đối cùng tỷ lệ phần trăm giúp tránh cách diễn giải sai khi số mẫu giữa hai lớp khác nhau.

\section{Kết quả trên bộ dữ liệu DroneDetect v2}

\subsection{Kết quả Linear Probing}

Trong thí nghiệm linear probing, ảnh spectrogram được đưa qua backbone đã đóng băng để trích xuất embedding. Các embedding của tập huấn luyện được dùng để huấn luyện SVM, trong khi embedding của tập kiểm thử chỉ được sử dụng để dự đoán. Kết quả cần được tổng hợp theo Bảng~\ref{tab:linear_probe_benchmark}.

\begin{table}[H]
    \centering
    \caption{Kết quả linear probing trên DroneDetect v2}
    \label{tab:linear_probe_benchmark}
    \renewcommand{\arraystretch}{1.3}
    \begin{tabular}{|p{0.25\textwidth}|c|c|c|c|}
        \hline
        \textbf{Backbone} & \textbf{Accuracy (\%)} & \textbf{Precision (\%)} & \textbf{Recall (\%)} & \textbf{F1-score (\%)} \\ \hline
        ResNet50 & -- & -- & -- & -- \\ \hline
        EfficientNet-B2 & -- & -- & -- & -- \\ \hline
        ConvNeXtV2 & -- & -- & -- & -- \\ \hline
        SigLIP & -- & -- & -- & -- \\ \hline
        DINOv2 & -- & -- & -- & -- \\ \hline
    \end{tabular}
\end{table}

Theo kết quả tổng hợp hiện có, accuracy của cả năm cấu hình trong Bảng~\ref{tab:linear_probe_benchmark} đều dưới $50\%$. Xu hướng này cho thấy embedding cố định chưa tạo ra không gian đặc trưng phù hợp để SVM phân tách ổn định hai lớp trong cấu hình thí nghiệm. Nếu tập kiểm thử cân bằng, các giá trị dưới $50\%$ còn thấp hơn mức đạt được khi luôn chọn một lớp; nếu tập không cân bằng, cần đối chiếu phân bố lớp trước khi so sánh với mức nền.

Nguyên nhân thứ nhất có thể xuất phát từ khoảng cách miền dữ liệu. Các backbone được tiền huấn luyện chủ yếu trên ảnh tự nhiên hoặc cặp ảnh--văn bản, trong khi spectrogram biểu diễn năng lượng theo thời gian và tần số. Các texture, cạnh và bố cục có ý nghĩa trong ảnh tự nhiên không nhất thiết tương ứng với dải truyền, burst hoặc nền nhiễu RF. Khi backbone bị đóng băng, các tầng không có cơ hội điều chỉnh bộ lọc và vùng chú ý theo cấu trúc này.

Nguyên nhân thứ hai liên quan đến cách tổng hợp embedding. Một vector toàn cục có thể làm giảm thông tin về vị trí tương đối của tín hiệu trên trục tần số hoặc thời gian. Nếu hai lớp khác nhau chủ yếu ở một vùng nhỏ của spectrogram, phép pooling toàn cục có thể pha loãng đặc trưng cần thiết. SVM chỉ nhận vector cuối nên không thể phục hồi thông tin đã mất ở bước trích xuất.

Nguyên nhân thứ ba có thể liên quan đến chuẩn hóa embedding và siêu tham số SVM. Tuy nhiên, vì cả năm backbone đều dưới $50\%$, việc chỉ điều chỉnh SVM khó giải thích toàn bộ xu hướng nếu quy trình đã sử dụng tập xác thực hợp lệ. Các giá trị precision, recall và ma trận nhầm lẫn sau khi được bổ sung sẽ cho biết mô hình đang thiên về \textit{Drone} hay \textit{Non-drone}, từ đó xác định lỗi chủ yếu đến từ ranh giới phân loại hay từ biểu diễn không phân tách.

\subsection{Kết quả Fine-tuning}

Trong thí nghiệm fine-tuning, lớp phân loại cuối được thay bằng lớp hai đầu ra và toàn bộ backbone được cập nhật trên dữ liệu huấn luyện. Checkpoint dùng để kiểm thử phải được lựa chọn bằng tập xác thực. Kết quả được trình bày theo Bảng~\ref{tab:finetune_benchmark}.

\begin{table}[H]
    \centering
    \caption{Kết quả fine-tuning trên DroneDetect v2}
    \label{tab:finetune_benchmark}
    \renewcommand{\arraystretch}{1.3}
    \begin{tabular}{|p{0.25\textwidth}|c|c|c|c|}
        \hline
        \textbf{Backbone} & \textbf{Accuracy (\%)} & \textbf{Precision (\%)} & \textbf{Recall (\%)} & \textbf{F1-score (\%)} \\ \hline
        ResNet50 & -- & -- & -- & -- \\ \hline
        EfficientNet-B2 & -- & -- & -- & -- \\ \hline
        ConvNeXtV2 & -- & -- & -- & -- \\ \hline
        SigLIP & -- & -- & -- & -- \\ \hline
        DINOv2 & -- & -- & -- & -- \\ \hline
    \end{tabular}
\end{table}

Sau khi điền số liệu, nhận xét giữa các backbone cần dựa trực tiếp vào chênh lệch trong Bảng~\ref{tab:finetune_benchmark}. Backbone có accuracy cao nhất chỉ được xem là tốt nhất về tỷ lệ dự đoán đúng tổng thể. Nếu một mô hình khác có recall cao hơn, mô hình đó bỏ sót ít drone hơn; nếu precision cao hơn, mô hình tạo ít cảnh báo sai hơn. Vì vậy, việc lựa chọn mô hình phải gắn với mục tiêu vận hành thay vì chỉ dựa trên thứ hạng accuracy.

Fine-tuning cho phép các tầng đầu và tầng sâu thay đổi theo spectrogram. Những bộ lọc ban đầu có thể điều chỉnh để phản ứng với cấu trúc năng lượng đơn sắc, còn các tầng sau có thể học quan hệ giữa vùng tần số và thời gian. Đây là cơ chế có khả năng giải thích mức cải thiện so với embedding cố định. Tuy nhiên, nếu một backbone đạt kết quả gần tuyệt đối trên benchmark, cần kiểm tra cách chia tập theo nguồn để loại trừ rò rỉ giữa các đoạn cùng bản ghi.

% Gợi ý chèn hình: Ma trận nhầm lẫn của năm mô hình fine-tuning trên tập kiểm thử DroneDetect v2.

\subsection{So sánh Linear Probing và Fine-tuning}

Bảng~\ref{tab:lp_ft_comparison} được sử dụng để so sánh trực tiếp hai chiến lược trên cùng backbone. Cột chênh lệch phải được tính từ các giá trị đã điền trong hai bảng trước, không nhập bằng ước lượng.

\begin{table}[H]
    \centering
    \caption{So sánh linear probing và fine-tuning trên DroneDetect v2}
    \label{tab:lp_ft_comparison}
    \renewcommand{\arraystretch}{1.3}
    \begin{tabular}{|p{0.24\textwidth}|c|c|c|c|}
        \hline
        \textbf{Backbone} & $\mathbf{Acc_{LP}}$ \textbf{(\%)} & $\mathbf{Acc_{FT}}$ \textbf{(\%)} & $\boldsymbol{\Delta Acc}$ & $\boldsymbol{\Delta F1}$ \\ \hline
        ResNet50 & -- & -- & [TÍNH] & [TÍNH] \\ \hline
        EfficientNet-B2 & -- & -- & [TÍNH] & [TÍNH] \\ \hline
        ConvNeXtV2 & -- & -- & [TÍNH] & [TÍNH] \\ \hline
        SigLIP & -- & -- & [TÍNH] & [TÍNH] \\ \hline
        DINOv2 & -- & -- & [TÍNH] & [TÍNH] \\ \hline
    \end{tabular}
\end{table}

% Gợi ý chèn hình: Biểu đồ cột nhóm so sánh accuracy của Linear Probing và Fine-tuning cho từng backbone.

Với thông tin hiện có rằng toàn bộ linear probing dưới $50\%$, fine-tuning chỉ được kết luận vượt trội cho từng backbone sau khi cột $\mathrm{Acc_{FT}}$ được điền và lớn hơn giá trị tương ứng. Nếu điều kiện này đúng cho cả năm dòng, kết quả hỗ trợ nhận định rằng thích nghi toàn bộ backbone hiệu quả hơn việc chỉ huấn luyện SVM trên embedding cố định trong bài toán này. Mức cải thiện tuyệt đối cần được báo cáo bằng điểm phần trăm, không dùng tỷ lệ phần trăm tương đối nếu không nêu rõ cách tính.

Nguyên nhân chính của chênh lệch dự kiến là fine-tuning tối ưu trực tiếp biểu diễn theo nhãn RF. Linear probing chỉ dịch chuyển ranh giới quyết định trong một không gian cố định, trong khi fine-tuning đồng thời thay đổi không gian đặc trưng và ranh giới phân loại. Điều này đặc biệt quan trọng khi tín hiệu drone được phân biệt bằng các cấu trúc không phải mục tiêu của quá trình tiền huấn luyện ban đầu.

Mặt khác, kết quả benchmark cao hơn không tự động chứng minh fine-tuning tổng quát hóa tốt hơn. Quá trình cập nhật toàn bộ tham số có thể làm mô hình học cả đặc trưng tín hiệu lẫn đặc trưng riêng của thiết bị và pipeline DroneDetect v2. Do đó, kết quả trên Drone1 và Drone2 được sử dụng như phép kiểm tra bổ sung cho lợi ích của fine-tuning.

\section{Kết quả trên dữ liệu RF thực tế}

\subsection{Kịch bản thử nghiệm}

Kịch bản thử nghiệm sử dụng các mô hình đã được huấn luyện và lựa chọn trên DroneDetect v2. Không thực hiện cập nhật tham số bằng Drone1 hoặc Drone2. Tín hiệu RF được thu bằng bladeRF tại băng 2,4~GHz, lưu dưới dạng I/Q, phân đoạn và chuyển thành spectrogram. Ảnh được chuẩn hóa theo cấu hình tương ứng của từng backbone trước khi đưa qua mô hình. Luồng xử lý có thể biểu diễn như sau:
\begin{equation}
    \text{RF thực tế}
    \longrightarrow
    \text{Spectrogram}
    \longrightarrow
    \text{Mô hình}
    \longrightarrow
    \text{Kết quả nhận dạng}.
\end{equation}

Mỗi tập Drone1 và Drone2 gồm 200 spectrogram mang nhãn \textit{Drone}. Mô hình đưa ra dự đoán nhị phân cho từng ảnh. Tỷ lệ nhận diện được tính bằng số ảnh dự đoán đúng là drone chia cho 200. Vì không có mẫu \textit{Non-drone}, thí nghiệm này tập trung vào khả năng phát hiện và số lượng bỏ sót, không đánh giá khả năng hạn chế cảnh báo giả.

% Gợi ý chèn hình: Pipeline suy luận trên dữ liệu thực tế từ bladeRF, tệp I/Q, spectrogram đến nhãn dự đoán.

\subsection{Thử nghiệm theo khoảng cách thu}

Bên cạnh hai tập Drone1 và Drone2, một thí nghiệm riêng được thiết kế để khảo sát ảnh hưởng của khoảng cách giữa nguồn phát và hệ thống thu. Tín hiệu được thu tại bốn khoảng cách 5~m, 10~m, 15~m và 20~m. Mục tiêu của thí nghiệm là xác định mức thay đổi của tỷ lệ nhận diện khi công suất tín hiệu tại đầu thu suy giảm theo khoảng cách, đồng thời đánh giá độ ổn định của từng backbone trong các điều kiện thu khác nhau.

Để khoảng cách là biến số chính, các yếu tố còn lại cần được giữ ổn định trong toàn bộ quá trình thu. Cùng một drone, bladeRF, ăng-ten, tần số trung tâm, tốc độ lấy mẫu, băng thông, độ lợi và thời lượng thu được sử dụng tại bốn vị trí. Hướng và độ cao tương đối của ăng-ten cũng cần được duy trì trong phạm vi có thể kiểm soát. Không nên điều chỉnh gain riêng cho từng khoảng cách nếu mục tiêu là quan sát trực tiếp ảnh hưởng của suy hao đường truyền. Nếu bắt buộc phải thay đổi gain để tránh mất tín hiệu, giá trị tương ứng phải được báo cáo và kết quả không còn phản ánh riêng tác động của khoảng cách.

Mỗi khoảng cách cần có cùng số lượng spectrogram và được tạo bằng cùng cấu hình phân đoạn, STFT, chuẩn hóa và resize. Các đoạn nên được lấy từ những phiên thu độc lập hoặc từ các khoảng thời gian không chồng lấn để giảm tương quan giữa các mẫu. Bảng~\ref{tab:distance_capture_setup} trình bày thông tin cần ghi nhận cho từng khoảng cách.

\begin{table}[H]
    \centering
    \caption{Cấu hình dữ liệu RF thu theo khoảng cách}
    \label{tab:distance_capture_setup}
    \renewcommand{\arraystretch}{1.3}
    \begin{tabular}{|c|c|c|c|c|p{0.22\textwidth}|}
        \hline
        \textbf{Khoảng cách} & \textbf{Số ảnh} & \textbf{Gain} & \textbf{SNR TB (dB)} & \textbf{Công suất TB (dB)} & \textbf{Điều kiện thu} \\ \hline
        5~m & -- & -- & -- & -- & -- \\ \hline
        10~m & -- & -- & -- & -- & -- \\ \hline
        15~m & -- & -- & -- & -- & -- \\ \hline
        20~m & -- & -- & -- & -- & -- \\ \hline
    \end{tabular}
\end{table}

Trong điều kiện truyền sóng lý tưởng, công suất thu có xu hướng giảm khi khoảng cách tăng. Tuy nhiên, môi trường thực tế còn chịu ảnh hưởng của phản xạ, che khuất, giao thoa đa đường và nguồn nhiễu trong băng 2,4~GHz. Vì vậy, khoảng cách không nên được dùng thay thế trực tiếp cho SNR. Việc báo cáo đồng thời khoảng cách, SNR và công suất trung bình giúp xác định hiệu năng giảm do tín hiệu yếu hơn hay do một điều kiện nhiễu cục bộ tại vị trí thu.

Các mô hình được suy luận trên từng tập khoảng cách mà không huấn luyện lại bằng dữ liệu này. Vì các tập chỉ chứa tín hiệu drone, chỉ số chính là tỷ lệ nhận diện, tương đương recall của lớp \textit{Drone}. Kết quả được trình bày trong Bảng~\ref{tab:distance_recognition_results}.

\begin{table}[H]
    \centering
    \caption{Tỷ lệ nhận diện drone theo khoảng cách thu}
    \label{tab:distance_recognition_results}
    \renewcommand{\arraystretch}{1.3}
    \begin{tabular}{|p{0.25\textwidth}|c|c|c|c|c|}
        \hline
        \textbf{Mô hình} & \textbf{5 m (\%)} & \textbf{10 m (\%)} & \textbf{15 m (\%)} & \textbf{20 m (\%)} & \textbf{Trung bình (\%)} \\ \hline
        ResNet50 & -- & -- & -- & -- & -- \\ \hline
        EfficientNet-B2 & -- & -- & -- & -- & -- \\ \hline
        ConvNeXtV2 & -- & -- & -- & -- & -- \\ \hline
        SigLIP & -- & -- & -- & -- & -- \\ \hline
        DINOv2 & -- & -- & -- & -- & -- \\ \hline
    \end{tabular}
\end{table}

% Gợi ý chèn hình: Biểu đồ đường biểu diễn tỷ lệ nhận diện của từng backbone tại các khoảng cách 5 m, 10 m, 15 m và 20 m.

Mức suy giảm từ khoảng cách gần nhất đến xa nhất của mô hình $m$ được tính bởi
\begin{equation}
    \Delta R_m
    =
    R_{m,5\mathrm{m}}
    -
    R_{m,20\mathrm{m}},
    \label{eq:distance_recall_drop}
\end{equation}
trong đó $R_{m,5\mathrm{m}}$ và $R_{m,20\mathrm{m}}$ lần lượt là tỷ lệ nhận diện tại 5~m và 20~m. Giá trị $\Delta R_m$ nhỏ cho thấy mô hình duy trì kết quả ổn định hơn trong dải khoảng cách khảo sát, nhưng chỉ có ý nghĩa tích cực khi tỷ lệ nhận diện tuyệt đối vẫn đủ cao.

\begin{table}[H]
    \centering
    \caption{Mức thay đổi tỷ lệ nhận diện theo khoảng cách}
    \label{tab:distance_degradation}
    \renewcommand{\arraystretch}{1.3}
    \begin{tabular}{|p{0.25\textwidth}|c|c|c|c|}
        \hline
        \textbf{Mô hình} & $\mathbf{R_{5m}}$ \textbf{(\%)} & $\mathbf{R_{20m}}$ \textbf{(\%)} & $\boldsymbol{\Delta R}$ \textbf{(điểm \%)} & \textbf{Xu hướng} \\ \hline
        ResNet50 & -- & -- & -- & -- \\ \hline
        EfficientNet-B2 & -- & -- & -- & -- \\ \hline
        ConvNeXtV2 & -- & -- & -- & -- \\ \hline
        SigLIP & -- & -- & -- & -- \\ \hline
        DINOv2 & -- & -- & -- & -- \\ \hline
    \end{tabular}
\end{table}

Kết quả không nhất thiết giảm đơn điệu từ 5~m đến 20~m. Một vị trí xa hơn vẫn có thể cho SNR tốt hơn nếu ít vật cản, ít nhiễu đồng kênh hoặc có điều kiện đa đường thuận lợi. Do đó, nhận xét phải dựa trên số liệu trong Bảng~\ref{tab:distance_capture_setup} và Bảng~\ref{tab:distance_recognition_results}. Nếu tỷ lệ nhận diện giảm đồng thời với SNR, kết quả hỗ trợ nhận định rằng suy hao tín hiệu theo khoảng cách làm giảm khả năng nhận dạng. Nếu hai đại lượng không biến đổi cùng chiều, cần xem xét thêm PSD, clipping, stripe artifact và điều kiện truyền sóng tại từng vị trí.

\subsection{Kết quả nhận dạng}

Kết quả nhận dạng của năm mô hình trên Drone1 và Drone2 được trình bày trong Bảng~\ref{tab:real_recognition_results}. Số mẫu nhận diện đúng được suy ra trực tiếp bằng tỷ lệ nhân với 200; số mẫu bỏ sót là phần còn lại.

\begin{table}[H]
    \centering
    \caption{Kết quả nhận diện drone trên hai tập dữ liệu RF tự thu}
    \label{tab:real_recognition_results}
    \renewcommand{\arraystretch}{1.3}
    \begin{tabular}{|p{0.22\textwidth}|c|c|c|c|c|}
        \hline
        \textbf{Mô hình} & \textbf{Drone1 (\%)} & \textbf{Đúng/200} & \textbf{Drone2 (\%)} & \textbf{Đúng/200} & \textbf{Trung bình (\%)} \\ \hline
        DINOv2 & 91,0 & 182 & 99,0 & 198 & 95,0 \\ \hline
        EfficientNet-B2 & 96,5 & 193 & 95,5 & 191 & 96,0 \\ \hline
        SigLIP & 84,5 & 169 & 96,5 & 193 & 90,5 \\ \hline
        ConvNeXtV2 & 49,5 & 99 & 61,0 & 122 & 55,25 \\ \hline
        ResNet50 & 20,5 & 41 & 16,5 & 33 & 18,5 \\ \hline
    \end{tabular}
\end{table}

EfficientNet-B2 có tỷ lệ trung bình cao nhất, đạt $96,0\%$, tương ứng nhận diện đúng tổng cộng 384 trên 400 ảnh. DINOv2 đứng sau với trung bình $95,0\%$ và 380 ảnh đúng. Chênh lệch trung bình giữa hai mô hình là 1 điểm phần trăm, tương ứng 4 ảnh trên toàn bộ dữ liệu thực tế. Vì quy mô đánh giá chỉ gồm 400 ảnh và chưa có khoảng tin cậy theo phiên thu, chênh lệch này cho thấy EfficientNet-B2 đạt kết quả quan sát cao hơn nhưng chưa đủ để khẳng định ưu thế tổng quát trong mọi điều kiện.

SigLIP đạt trung bình $90,5\%$, tương ứng 362 trên 400 ảnh đúng. Kết quả này thấp hơn EfficientNet-B2 5,5 điểm phần trăm và thấp hơn DINOv2 4,5 điểm phần trăm. Tuy nhiên, SigLIP vẫn nhận diện đúng hơn chín phần mười tổng số mẫu, cho thấy backbone có khả năng chuyển giao đáng kể sang hai tập RF tự thu.

ConvNeXtV2 đạt trung bình $55,25\%$, tương ứng 221 trên 400 ảnh đúng. Giá trị này thấp hơn SigLIP 35,25 điểm phần trăm và chỉ cao hơn một mức nhận diện $50\%$ là 5,25 điểm phần trăm. ResNet50 đạt trung bình $18,5\%$, tương ứng 74 trên 400 ảnh đúng và bỏ sót 326 ảnh. Số liệu cho thấy hai backbone này tổng quát hóa kém hơn rõ rệt trong kịch bản tự thu, dù nguyên nhân cần được đối chiếu với kết quả benchmark và ma trận nhầm lẫn.

Kết quả cũng cho thấy không thể kết luận Transformer luôn tốt hơn CNN. EfficientNet-B2 là một CNN nhưng đạt trung bình $96,0\%$, cao hơn DINOv2 và SigLIP trong phép đo này. Đồng thời, ConvNeXtV2 và ResNet50 cũng là CNN nhưng cho kết quả thấp hơn đáng kể. Sự khác biệt giữa ba CNN chứng minh khả năng tổng quát hóa phụ thuộc vào từng backbone, trọng số tiền huấn luyện và quá trình fine-tuning, không chỉ phụ thuộc tên nhóm kiến trúc.

% Gợi ý chèn hình: Biểu đồ cột nhóm thể hiện tỷ lệ nhận diện của năm mô hình trên Drone1 và Drone2.

\subsection{Phân tích mức suy giảm hiệu năng}

Mức suy giảm từ DroneDetect v2 sang dữ liệu thực tế chỉ có thể tính chính xác sau khi điền kết quả benchmark. Bảng~\ref{tab:benchmark_real_drop} tạo khung đối chiếu giữa recall lớp drone trên benchmark và tỷ lệ nhận diện trung bình trên Drone1, Drone2. Recall được sử dụng thay cho accuracy để bảo đảm hai đại lượng cùng đo khả năng nhận đúng lớp drone.

\begin{table}[H]
    \centering
    \caption{Đối chiếu khả năng nhận diện drone giữa benchmark và dữ liệu tự thu}
    \label{tab:benchmark_real_drop}
    \renewcommand{\arraystretch}{1.3}
    \begin{tabular}{|p{0.24\textwidth}|c|c|c|}
        \hline
        \textbf{Mô hình} & \textbf{Recall benchmark (\%)} & \textbf{Recall thực tế TB (\%)} & \textbf{Chênh lệch (điểm \%)} \\ \hline
        DINOv2 & -- & 95,0 & [TÍNH] \\ \hline
        EfficientNet-B2 & -- & 96,0 & [TÍNH] \\ \hline
        SigLIP & -- & 90,5 & [TÍNH] \\ \hline
        ConvNeXtV2 & -- & 55,25 & [TÍNH] \\ \hline
        ResNet50 & -- & 18,5 & [TÍNH] \\ \hline
    \end{tabular}
\end{table}

Domain shift là nguyên nhân chính cần xem xét khi benchmark và dữ liệu thực tế có chênh lệch. DroneDetect v2 và dữ liệu bladeRF có thể khác nhau về đáp ứng tần số của thiết bị, mức khuếch đại, băng thông, tốc độ lấy mẫu và nhiễu nền. Ngay cả khi cùng chứa tín hiệu drone, các khác biệt này làm thay đổi độ tương phản, độ rộng dải và texture của spectrogram.

Môi trường 2,4~GHz còn chứa các nguồn phát khác như Wi-Fi và Bluetooth. Năng lượng từ các nguồn này có thể chồng lấn với vùng tần số của tín hiệu cần phát hiện hoặc làm nền phổ không giống dữ liệu benchmark. Nếu mô hình học dựa nhiều vào mức nền, vị trí dải hoặc artifact của nguồn dữ liệu huấn luyện, hiệu năng sẽ giảm khi các dấu hiệu đó thay đổi.

Kết quả thực tế trong Bảng~\ref{tab:real_recognition_results} cho thấy mức ảnh hưởng của domain shift không đồng đều giữa các backbone. EfficientNet-B2 và DINOv2 đều duy trì trung bình ít nhất $95\%$, trong khi ResNet50 chỉ đạt $18,5\%$. Do các mô hình nhận cùng các tập Drone1 và Drone2, chênh lệch này cho thấy biểu diễn và quá trình thích nghi của từng backbone có vai trò đáng kể. Tuy nhiên, chưa thể tách riêng ảnh hưởng của kiến trúc, trọng số tiền huấn luyện và siêu tham số nếu chưa có thí nghiệm kiểm soát từng yếu tố.

\subsection{So sánh khả năng tổng quát hóa}

DINOv2 đạt $91,0\%$ trên Drone1 và $99,0\%$ trên Drone2. Chênh lệch 8 điểm phần trăm tương ứng với 16 ảnh. Mô hình nhận diện tốt hơn trên Drone2, nơi chỉ có 2 ảnh bị bỏ sót, so với 18 ảnh trên Drone1. Điều này cho thấy DINOv2 có kết quả tổng thể cao nhưng vẫn nhạy với khác biệt giữa hai tập.

EfficientNet-B2 đạt $96,5\%$ trên Drone1 và $95,5\%$ trên Drone2. Chênh lệch chỉ 1 điểm phần trăm, tương ứng 2 ảnh. Đây là mô hình có kết quả ổn định nhất giữa hai tập trong năm backbone, đồng thời có tỷ lệ trung bình cao nhất. Số liệu này cho thấy EfficientNet-B2 là lựa chọn nổi bật cho kịch bản tự thu đã thực hiện.

SigLIP đạt $84,5\%$ trên Drone1 và $96,5\%$ trên Drone2. Chênh lệch 12 điểm phần trăm tương ứng 24 ảnh, lớn nhất trong năm mô hình. Kết quả cho thấy SigLIP có khả năng nhận diện cao trên Drone2 nhưng giảm đáng kể trên Drone1. Vì vậy, giá trị trung bình $90,5\%$ cần được diễn giải cùng độ biến động giữa hai tập.

ConvNeXtV2 tăng từ $49,5\%$ trên Drone1 lên $61,0\%$ trên Drone2, chênh lệch 11,5 điểm phần trăm. Mô hình nhận diện đúng 99 ảnh ở Drone1 và 122 ảnh ở Drone2. Mặc dù kết quả Drone2 cao hơn, cả hai giá trị đều thấp hơn đáng kể so với ba mô hình dẫn đầu.

ResNet50 đạt $20,5\%$ trên Drone1 và $16,5\%$ trên Drone2. Mô hình chỉ nhận diện đúng lần lượt 41 và 33 ảnh, đồng thời có chênh lệch 4 điểm phần trăm. Độ chênh nhỏ không thể được xem là dấu hiệu ổn định tích cực khi mức hiệu năng tuyệt đối rất thấp. Trong trường hợp này, mô hình ổn định ở mức không đáp ứng yêu cầu phát hiện.

\begin{table}[H]
    \centering
    \caption{Độ ổn định giữa Drone1 và Drone2}
    \label{tab:real_stability}
    \renewcommand{\arraystretch}{1.3}
    \begin{tabular}{|p{0.27\textwidth}|c|c|c|}
        \hline
        \textbf{Mô hình} & \textbf{Trung bình (\%)} & \textbf{Chênh lệch tuyệt đối (\%)} & \textbf{Tổng số bỏ sót} \\ \hline
        DINOv2 & 95,0 & 8,0 & 20 \\ \hline
        EfficientNet-B2 & 96,0 & 1,0 & 16 \\ \hline
        SigLIP & 90,5 & 12,0 & 38 \\ \hline
        ConvNeXtV2 & 55,25 & 11,5 & 179 \\ \hline
        ResNet50 & 18,5 & 4,0 & 326 \\ \hline
    \end{tabular}
\end{table}

Bảng~\ref{tab:real_stability} củng cố nhận xét rằng EfficientNet-B2 đạt sự kết hợp tốt nhất giữa tỷ lệ trung bình và độ ổn định trên hai tập. DINOv2 có trung bình thấp hơn 1 điểm phần trăm nhưng bỏ sót nhiều hơn 4 ảnh. SigLIP đứng thứ ba về trung bình và có mức chênh lệch lớn nhất. ConvNeXtV2 và ResNet50 có số lượng bỏ sót lần lượt là 179 và 326, cho thấy hạn chế rõ rệt trong kịch bản thực tế.

\section{Phân tích ảnh hưởng của tham số tạo spectrogram}

Nếu tỷ lệ nhận diện tăng theo SNR trong Bảng~\ref{tab:snr_analysis}, số liệu sẽ hỗ trợ mối liên hệ giữa độ rõ của tín hiệu và khả năng phân loại. Nếu xu hướng không đơn điệu, nguyên nhân có thể là SNR không phản ánh đầy đủ loại nhiễu, vị trí tín hiệu hoặc artifact. Do đó, SNR cần được phân tích cùng PSD và hình dạng spectrogram.

\subsection{Thiết kế khảo sát tham số STFT}

Spectrogram không chỉ phụ thuộc vào chất lượng tín hiệu I/Q mà còn phụ thuộc vào cấu hình STFT. Kích thước cửa sổ, kích thước FFT, hop length và loại cửa sổ quyết định cách năng lượng được phân bố trên hai trục thời gian--tần số. Một cấu hình không phù hợp có thể làm mất các burst ngắn, làm mờ biên dải phổ hoặc tạo ảnh có số hàng và số cột mất cân đối trước bước resize. Vì vậy, ảnh hưởng của tham số STFT cần được khảo sát như một phần của thực nghiệm thay vì xem đây là một bước tiền xử lý cố định.

Khảo sát được thực hiện theo nguyên tắc chỉ thay đổi một yếu tố tại mỗi nhóm thí nghiệm. Một cấu hình cơ sở được lựa chọn gồm độ dài cửa sổ $N_w$, kích thước FFT $N_{\mathrm{FFT}}$, hop length $H$ và cửa sổ Hann. Khi khảo sát $N_{\mathrm{FFT}}$, các tham số còn lại được giữ nguyên; quy tắc tương tự được áp dụng cho $N_w$, $H$ và loại cửa sổ. Toàn bộ spectrogram của tập huấn luyện, xác thực và kiểm thử phải được sinh lại theo cấu hình đang xét.

Mỗi cấu hình STFT cần được huấn luyện và đánh giá độc lập bằng cùng một backbone, cùng cách chia dữ liệu, cùng số epoch và cùng quy tắc lựa chọn checkpoint. Việc chỉ thay đổi STFT ở tập kiểm thử trong khi giữ mô hình huấn luyện bằng cấu hình khác sẽ đo độ nhạy trước sai lệch tiền xử lý, không đo chất lượng vốn có của cấu hình STFT. Do chi phí tính toán lớn, thí nghiệm lựa chọn tham số có thể sử dụng một backbone đại diện, ưu tiên EfficientNet-B2 hoặc DINOv2 dựa trên kết quả thực tế, sau đó xác nhận cấu hình được chọn trên các backbone còn lại.

\begin{table}[H]
    \centering
    \caption{Các nhóm cấu hình dùng để khảo sát tham số STFT}
    \label{tab:stft_experiment_design}
    \renewcommand{\arraystretch}{1.3}
    \begin{tabular}{|p{0.23\textwidth}|p{0.27\textwidth}|p{0.39\textwidth}|}
        \hline
        \textbf{Tham số khảo sát} & \textbf{Các giá trị} & \textbf{Các tham số giữ cố định} \\ \hline
        Độ dài cửa sổ $N_w$ & -- & $N_{\mathrm{FFT}}$, $H$, loại cửa sổ và quy tắc chuẩn hóa \\ \hline
        Kích thước FFT $N_{\mathrm{FFT}}$ & -- & $N_w$, $H$, loại cửa sổ và quy tắc chuẩn hóa \\ \hline
        Hop length $H$ & -- & $N_w$, $N_{\mathrm{FFT}}$, loại cửa sổ và quy tắc chuẩn hóa \\ \hline
        Loại cửa sổ & Hann, Hamming, Blackman -- & $N_w$, $N_{\mathrm{FFT}}$, $H$ và quy tắc chuẩn hóa \\ \hline
    \end{tabular}
\end{table}

Các giá trị trong Bảng~\ref{tab:stft_experiment_design} chỉ là khung tham khảo và cần được điều chỉnh theo tốc độ lấy mẫu, độ dài đoạn I/Q và cấu hình đã sử dụng. Không nên khảo sát các giá trị tạo ra quá ít khung thời gian hoặc kích thước FFT nhỏ hơn độ dài cửa sổ. Khi $N_{\mathrm{FFT}}>N_w$, phần tăng độ phân giải lưới tần số do zero-padding không đồng nghĩa với tăng độ phân giải vật lý của tín hiệu, do đó kết quả cần được diễn giải thận trọng.

\subsection{Ảnh hưởng của độ dài cửa sổ và kích thước FFT}

Độ dài cửa sổ chi phối sự đánh đổi giữa độ phân giải thời gian và độ phân giải tần số. Cửa sổ dài quan sát nhiều mẫu hơn trong mỗi phép FFT, nhờ đó có khả năng phân biệt các thành phần tần số gần nhau tốt hơn. Tuy nhiên, năng lượng của các sự kiện ngắn bị tổng hợp trong một khoảng thời gian dài hơn, làm giảm khả năng định vị burst. Cửa sổ ngắn bảo toàn biến thiên nhanh theo thời gian nhưng làm các dải phổ rộng hơn và khó tách các thành phần gần nhau.

Kích thước FFT quyết định số bin trên trục tần số. Khi $N_{\mathrm{FFT}}$ bằng độ dài cửa sổ, tăng kích thước FFT đồng thời với tăng cửa sổ có thể cải thiện độ phân giải tần số nhưng làm giảm độ phân giải thời gian. Khi chỉ tăng $N_{\mathrm{FFT}}$ và giữ nguyên $N_w$, tín hiệu được zero-padding trước FFT; ảnh có nhiều bin hơn nhưng không chứa thêm thông tin quan sát. Sau bước resize về kích thước đầu vào cố định, số bin quá lớn còn có thể bị nội suy và không tạo ra lợi ích tương ứng cho mô hình.

\begin{table}[H]
    \centering
    \caption{Kết quả khảo sát độ dài cửa sổ và kích thước FFT}
    \label{tab:stft_window_fft_results}
    \renewcommand{\arraystretch}{1.3}
    \begin{tabular}{|c|c|c|c|c|c|}
        \hline
        $\mathbf{N_w}$ & $\mathbf{N_{\mathrm{FFT}}}$ & \textbf{Số khung} & \textbf{Accuracy (\%)} & \textbf{Recall (\%)} & \textbf{F1-score (\%)} \\ \hline
        -- & -- & -- & -- & -- & -- \\ \hline
        -- & -- & -- & -- & -- & -- \\ \hline
        -- & -- & -- & -- & -- & -- \\ \hline
        -- & -- & -- & -- & -- & -- \\ \hline
    \end{tabular}
\end{table}

Nhận xét từ Bảng~\ref{tab:stft_window_fft_results} cần dựa trên chênh lệch số liệu và đặc điểm ảnh tương ứng. Nếu cấu hình có cửa sổ lớn làm recall giảm, một nguyên nhân có thể là các burst ngắn bị làm mờ theo thời gian. Nếu kích thước FFT nhỏ làm precision giảm, các dải phổ của drone và nhiễu có thể bị trộn trên lưới tần số thưa. Cấu hình tốt nhất không nhất thiết là cấu hình có $N_{\mathrm{FFT}}$ lớn nhất mà là cấu hình duy trì đủ thông tin trên cả hai trục sau khi resize.

% Gợi ý chèn hình: Cùng một đoạn I/Q được biểu diễn bằng nhiều cặp độ dài cửa sổ và kích thước FFT để minh họa sự đánh đổi thời gian--tần số.

\subsection{Ảnh hưởng của hop length và độ chồng lấn}

Hop length xác định khoảng dịch giữa hai cửa sổ STFT liên tiếp. Với độ dài cửa sổ $N_w$, tỷ lệ chồng lấn được tính bởi
\begin{equation}
    r_{\mathrm{overlap}}
    =
    1-\frac{H}{N_w}.
    \label{eq:stft_overlap}
\end{equation}
Hop length nhỏ tạo độ chồng lấn lớn, tăng số cột thời gian và giúp mô tả liên tục hơn các burst. Đổi lại, các cột kề nhau có tương quan cao, kích thước dữ liệu tăng và chi phí tạo spectrogram lớn hơn. Hop length lớn làm giảm chi phí nhưng có thể bỏ qua biến thiên xảy ra giữa hai vị trí cửa sổ.

\begin{table}[H]
    \centering
    \caption{Kết quả khảo sát hop length}
    \label{tab:stft_hop_results}
    \renewcommand{\arraystretch}{1.3}
    \begin{tabular}{|c|c|c|c|c|c|}
        \hline
        $\mathbf{H}$ & \textbf{Chồng lấn (\%)} & \textbf{Số khung} & \textbf{Accuracy (\%)} & \textbf{Recall (\%)} & \textbf{F1-score (\%)} \\ \hline
        -- & -- & -- & -- & -- & -- \\ \hline
        -- & -- & -- & -- & -- & -- \\ \hline
        -- & -- & -- & -- & -- & -- \\ \hline
    \end{tabular}
\end{table}

Nếu giảm hop length chỉ làm tăng số khung nhưng các chỉ số gần như không thay đổi, mức chồng lấn cao hơn không cung cấp thêm thông tin hữu ích sau resize. Nếu recall tăng, mô hình có thể được hưởng lợi từ việc các burst ngắn được biểu diễn liên tục hơn. Phân tích cần báo cáo thêm thời gian sinh dữ liệu hoặc kích thước lưu trữ để tránh lựa chọn một cấu hình có cải thiện rất nhỏ nhưng chi phí tăng lớn.

\subsection{Ảnh hưởng của loại cửa sổ và chuẩn hóa phổ}

Loại cửa sổ quyết định mức rò rỉ phổ và độ rộng búp chính. Cửa sổ Hann tạo sự cân bằng phổ biến giữa hai yếu tố này. Hamming có mức suy giảm búp phụ khác Hann, còn Blackman giảm rò rỉ phổ mạnh hơn nhưng làm búp chính rộng hơn. Sự khác biệt có thể ảnh hưởng đến độ sắc nét của dải tín hiệu và mức năng lượng nền trên spectrogram.

Ngoài loại cửa sổ, quy tắc chuyển công suất sang thang logarit và chuẩn hóa cường độ cũng ảnh hưởng trực tiếp đến độ tương phản ảnh. Chuẩn hóa riêng từng ảnh làm nổi bật cấu trúc yếu nhưng loại bỏ chênh lệch công suất tuyệt đối giữa các mẫu. Chuẩn hóa toàn cục giữ quan hệ mức năng lượng giữa các ảnh nhưng có thể làm tín hiệu yếu kém nổi bật. Hai yếu tố này nên được khảo sát sau khi đã cố định $N_w$, $N_{\mathrm{FFT}}$ và $H$.

\begin{table}[H]
    \centering
    \caption{Kết quả khảo sát cửa sổ và phương pháp chuẩn hóa spectrogram}
    \label{tab:stft_window_normalization_results}
    \renewcommand{\arraystretch}{1.3}
    \begin{tabular}{|p{0.20\textwidth}|p{0.25\textwidth}|c|c|c|}
        \hline
        \textbf{Cửa sổ} & \textbf{Chuẩn hóa} & \textbf{Accuracy (\%)} & \textbf{Recall (\%)} & \textbf{F1-score (\%)} \\ \hline
        Hann & -- & -- & -- & -- \\ \hline
        Hamming & -- & -- & -- & -- \\ \hline
        Blackman & -- & -- & -- & -- \\ \hline
        -- & -- & -- & -- & -- \\ \hline
    \end{tabular}
\end{table}

\subsection{Độ nhạy khi cấu hình STFT thay đổi giữa huấn luyện và suy luận}

Một thí nghiệm bổ sung là giữ nguyên mô hình được huấn luyện bằng cấu hình STFT cơ sở, sau đó thay đổi từng tham số khi tạo spectrogram cho tập kiểm thử. Thí nghiệm này mô phỏng tình huống pipeline triển khai không hoàn toàn trùng với pipeline huấn luyện. Mức suy giảm phản ánh độ nhạy của mô hình trước sai lệch tiền xử lý, không được dùng để kết luận cấu hình mới kém hơn về bản chất.

\begin{table}[H]
    \centering
    \caption{Độ nhạy trước sai lệch cấu hình STFT tại giai đoạn suy luận}
    \label{tab:stft_mismatch_results}
    \renewcommand{\arraystretch}{1.3}
    \begin{tabular}{|p{0.31\textwidth}|c|c|c|}
        \hline
        \textbf{Cấu hình suy luận} & \textbf{Accuracy (\%)} & \textbf{Recall (\%)} & \textbf{Mức giảm F1 (điểm \%)} \\ \hline
        Trùng cấu hình huấn luyện & -- & -- & 0 \\ \hline
        Thay đổi $N_{\mathrm{FFT}}$ & -- & -- & [TÍNH] \\ \hline
        Thay đổi hop length & -- & -- & [TÍNH] \\ \hline
        Thay đổi loại cửa sổ & -- & -- & [TÍNH] \\ \hline
        Thay đổi chuẩn hóa & -- & -- & [TÍNH] \\ \hline
    \end{tabular}
\end{table}

Nếu hiệu năng giảm mạnh khi chỉ thay đổi tham số STFT ở giai đoạn suy luận, mô hình đã phụ thuộc đáng kể vào hình học và phân bố cường độ của spectrogram trong tập huấn luyện. Kết quả này nhấn mạnh yêu cầu lưu cấu hình tiền xử lý cùng checkpoint và áp dụng chính xác cấu hình đó khi xử lý dữ liệu bladeRF.

\subsection{Tác động tới kết quả dự đoán}

Cùng một mô hình có thể cho kết quả khác nhau khi chất lượng spectrogram thay đổi vì biểu diễn đầu vào quyết định trực tiếp các kích hoạt trong backbone. Khi tín hiệu rõ, các vùng năng lượng có hình dạng và độ tương phản gần dữ liệu huấn luyện, nên vector đặc trưng có khả năng nằm gần các mẫu drone đã biết. Khi nhiễu, clipping hoặc artifact xuất hiện, vector đặc trưng có thể dịch chuyển qua ranh giới quyết định.

Chênh lệch giữa Drone1 và Drone2 trong Bảng~\ref{tab:real_recognition_results} cung cấp dấu hiệu rằng điều kiện dữ liệu ảnh hưởng khác nhau đến từng backbone. DINOv2 thay đổi 8 điểm phần trăm, SigLIP thay đổi 12 điểm và ConvNeXtV2 thay đổi 11,5 điểm, trong khi EfficientNet-B2 chỉ thay đổi 1 điểm. Tuy nhiên, chưa thể quy toàn bộ chênh lệch này cho chất lượng spectrogram nếu chưa đối chiếu SNR và artifact của hai tập. Điều kiện thu, loại tín hiệu và phân bố thời gian cũng có thể góp phần.

% Gợi ý chèn hình: Các mẫu dự đoán đúng và sai của từng mô hình, kèm nhãn, confidence, SNR và trạng thái clipping/artifact.

Phân tích lỗi nên ưu tiên các mẫu mà mô hình có độ tin cậy cao nhưng dự đoán sai. Những trường hợp này cho thấy mô hình không chỉ thiếu chắc chắn mà có thể đã học một dấu hiệu gây nhầm lẫn. Việc so sánh cùng một mẫu qua năm backbone cũng giúp xác định artifact ảnh hưởng chung đến pipeline hay chỉ ảnh hưởng mạnh đến một kiểu biểu diễn.

\section{Thảo luận}

\subsection{Ảnh hưởng của chiến lược huấn luyện}

Theo thông tin kết quả hiện có, accuracy của toàn bộ năm backbone trong chiến lược linear probing đều dưới $50\%$. Kết quả này cho thấy việc sử dụng trực tiếp embedding cố định kết hợp với SVM chưa phù hợp với cấu hình dữ liệu và tiền xử lý của đề tài. Không gian đặc trưng tiền huấn luyện chưa tạo được sự phân tách ổn định giữa \textit{Drone} và \textit{Non-drone}.

Khoảng cách miền giữa ảnh tiền huấn luyện và spectrogram là một nguyên nhân có thể giải thích xu hướng trên. Khi backbone bị đóng băng, các tầng vẫn ưu tiên những cấu trúc từng xuất hiện trong dữ liệu nguồn. Phép pooling thành embedding toàn cục cũng có thể làm giảm thông tin vị trí của dải phổ hoặc burst ngắn. SVM chỉ tối ưu ranh giới trên vector cuối nên không thể điều chỉnh các đặc trưng đã được trích xuất.

Bảng~\ref{tab:lp_ft_comparison} là căn cứ định lượng để đánh giá lợi ích của fine-tuning. Khi $\Delta \mathrm{Acc}$ và $\Delta \mathrm{F1}$ dương, kết quả chứng tỏ việc cập nhật toàn bộ backbone giúp mô hình thích nghi tốt hơn với cấu trúc thời gian--tần số. Fine-tuning đồng thời thay đổi không gian đặc trưng và lớp phân loại, thay vì chỉ thay đổi siêu phẳng quyết định như linear probing.

Tuy nhiên, fine-tuning có chi phí tính toán cao hơn và có thể học các đặc điểm riêng của DroneDetect v2. Vì vậy, cải thiện trên benchmark cần được xem xét cùng kết quả Drone1 và Drone2. Việc một số backbone đạt kết quả thực tế cao trong khi các backbone khác suy giảm mạnh cho thấy fine-tuning không tự động bảo đảm khả năng tổng quát hóa.

\subsection{Khả năng tổng quát hóa của các backbone}

Trên hai tập tự thu, EfficientNet-B2 có tỷ lệ trung bình cao nhất là $96,0\%$ và độ chênh giữa hai tập thấp nhất là 1 điểm phần trăm. DINOv2 đạt trung bình $95,0\%$, thấp hơn 1 điểm phần trăm. SigLIP đứng thứ ba với $90,5\%$. Ba mô hình này cho kết quả thực tế tốt nhất trong phạm vi năm backbone được khảo sát.

ConvNeXtV2 và ResNet50 chỉ đạt trung bình lần lượt $55,25\%$ và $18,5\%$. Chênh lệch lớn giữa các backbone trên cùng dữ liệu cho thấy khả năng tổng quát hóa phụ thuộc vào biểu diễn được học và quá trình tối ưu của từng mô hình. Độ ổn định cũng phải được xem xét cùng mức hiệu năng tuyệt đối. ResNet50 chỉ chênh 4 điểm phần trăm giữa Drone1 và Drone2, nhưng mức trung bình thấp nên không thể xem đây là sự ổn định có ý nghĩa tích cực.

Kết quả không cung cấp căn cứ để kết luận Transformer luôn tốt hơn CNN. EfficientNet-B2 là một CNN nhưng đạt tỷ lệ trung bình cao nhất. Trong khi đó, ConvNeXtV2 và ResNet50 cũng thuộc nhóm CNN nhưng có kết quả thấp hơn rõ rệt. Vì các mô hình còn khác nhau về dữ liệu tiền huấn luyện, quy mô và cách tổ chức đặc trưng, nhận xét phù hợp chỉ nên được đưa ra ở mức từng backbone.

\subsection{Sự khác biệt giữa benchmark và dữ liệu thực tế}

DroneDetect v2 tạo điều kiện huấn luyện và so sánh các mô hình trong một môi trường dữ liệu có kiểm soát. Tuy nhiên, dữ liệu bladeRF khác về phần cứng thu, đáp ứng tần số, mức khuếch đại, nhiễu nền và các nguồn phát cùng băng tần. Do đó, kết quả benchmark không thể thay thế hoàn toàn phép đánh giá trên dữ liệu thực tế.

Sự phân hóa từ $96,0\%$ của EfficientNet-B2 đến $18,5\%$ của ResNet50 cho thấy dữ liệu tự thu cung cấp thêm thông tin quan trọng về khả năng chuyển miền. Sau khi recall benchmark được bổ sung vào Bảng~\ref{tab:benchmark_real_drop}, mức chênh lệch sẽ cho biết mô hình nào suy giảm nhiều nhất. Nếu thứ hạng benchmark không trùng với thứ hạng thực tế, việc lựa chọn mô hình triển khai không nên chỉ dựa trên accuracy của DroneDetect v2.

Benchmark vẫn có vai trò cần thiết trong huấn luyện, lựa chọn checkpoint và so sánh có kiểm soát. Dữ liệu thực tế đóng vai trò bổ sung để kiểm tra độ bền vững trước domain shift. Hai phép đánh giá cần được sử dụng đồng thời thay vì xem một nguồn dữ liệu là sự thay thế cho nguồn còn lại.

\subsection{Ảnh hưởng của khoảng cách thu}

Thí nghiệm tại 5~m, 10~m, 15~m và 20~m bổ sung một yếu tố triển khai quan trọng cho đánh giá thực tế. Khoảng cách tăng thường làm giảm công suất tín hiệu tại đầu thu, từ đó giảm độ tương phản giữa cấu trúc drone và nền trên spectrogram. Tuy nhiên, khoảng cách chỉ là một biến đại diện; SNR và điều kiện truyền sóng mới phản ánh trực tiếp chất lượng tín hiệu mà mô hình nhận được.

Bảng~\ref{tab:distance_recognition_results} cho phép so sánh tỷ lệ nhận diện của từng backbone tại cùng một khoảng cách, trong khi Bảng~\ref{tab:distance_degradation} thể hiện độ ổn định từ 5~m đến 20~m. Mô hình phù hợp hơn cho triển khai cần duy trì cả tỷ lệ nhận diện cao và mức suy giảm thấp. Một mô hình có $\Delta R$ nhỏ nhưng hiệu năng thấp tại mọi khoảng cách không thể được xem là có khả năng hoạt động tốt.

Kết quả theo khoảng cách cần được đối chiếu với SNR trung bình trong Bảng~\ref{tab:distance_capture_setup}. Nếu hai backbone suy giảm khác nhau trên cùng dữ liệu, chênh lệch phản ánh độ bền vững khác nhau của biểu diễn. Nếu toàn bộ mô hình cùng giảm mạnh tại một vị trí, nguyên nhân có khả năng liên quan đến chất lượng phiên thu hoặc môi trường RF hơn là một kiến trúc riêng lẻ.

\subsection{Vai trò của cấu hình STFT}

Cấu hình STFT là một thành phần của mô hình thực nghiệm vì nó quyết định thông tin được biểu diễn trước khi ảnh đi vào backbone. Độ dài cửa sổ và kích thước FFT chi phối sự đánh đổi giữa độ phân giải thời gian và tần số, trong khi hop length quyết định mật độ quan sát theo thời gian. Loại cửa sổ và phương pháp chuẩn hóa tiếp tục làm thay đổi rò rỉ phổ và độ tương phản của ảnh.

Kết quả từ Bảng~\ref{tab:stft_window_fft_results}, Bảng~\ref{tab:stft_hop_results} và Bảng~\ref{tab:stft_window_normalization_results} cần được sử dụng để lựa chọn cấu hình dựa trên chỉ số phân loại thay vì chỉ dựa trên hình ảnh trực quan. Nếu nhiều cấu hình có hiệu năng gần nhau, cấu hình có chi phí tạo ảnh thấp hơn và ổn định hơn trên dữ liệu thực tế nên được ưu tiên.

Bảng~\ref{tab:stft_mismatch_results} đánh giá một khía cạnh khác là độ nhạy trước sai lệch tiền xử lý. Mức suy giảm lớn khi cấu hình suy luận khác cấu hình huấn luyện cho thấy tham số STFT phải được xem là một phần của checkpoint triển khai. Mô hình và pipeline tạo spectrogram cần được lưu, kiểm thử và phiên bản hóa cùng nhau.

\subsection{Giới hạn của thực nghiệm}

Drone1 và Drone2 chỉ chứa mẫu thuộc lớp \textit{Drone}. Vì vậy, kết quả trên hai tập phản ánh recall của lớp drone nhưng không đo được precision, specificity hoặc tỷ lệ cảnh báo giả. Một mô hình nhận diện đúng phần lớn mẫu drone vẫn có thể dự đoán nhầm nhiều tín hiệu nền thành drone nếu được triển khai trong môi trường mở.

Mỗi tập thực tế gồm 200 spectrogram và được tạo từ số lượng điều kiện thu còn giới hạn. Thí nghiệm tại bốn khoảng cách mở rộng phạm vi đánh giá nhưng vẫn chưa đại diện cho mọi khoảng cách, loại drone, cấu hình ăng-ten hoặc môi trường nhiễu. Các ảnh cắt từ cùng một bản ghi còn có thể tương quan với nhau, vì vậy số lượng ảnh không nhất thiết tương đương cùng số lượng quan sát độc lập.

Ảnh hưởng của SNR, clipping, stripe artifact và tham số STFT mới chỉ có thể kết luận định lượng sau khi hoàn thiện các bảng thí nghiệm tương ứng. Chênh lệch giữa Drone1 và Drone2 là dấu hiệu về tác động của điều kiện dữ liệu, nhưng chưa đủ để quy nguyên nhân cho một yếu tố riêng lẻ. Khi khảo sát STFT, chi phí tính toán có thể giới hạn số backbone được huấn luyện lại cho từng cấu hình; do đó, kết quả trên một backbone đại diện không nên được mặc nhiên suy rộng cho toàn bộ mô hình. Các giới hạn này cần được xem xét khi diễn giải kết quả và đề xuất phạm vi áp dụng của hệ thống.

\section{Kết chương}

Chương này đã trình bày thiết kế và kết quả thực nghiệm dùng để đánh giá chiến lược huấn luyện, backbone và khả năng tổng quát hóa của hệ thống phát hiện drone. Các thí nghiệm sử dụng DroneDetect v2 cho đánh giá benchmark và hai tập Drone1, Drone2 thu bằng bladeRF cho đánh giá thực tế. Accuracy, precision, recall, F1-score và ma trận nhầm lẫn được lựa chọn để mô tả các khía cạnh khác nhau của lỗi phân loại.

Kết quả hiện có cho thấy linear probing với SVM chưa phù hợp trong cấu hình thí nghiệm vì toàn bộ backbone đều đạt accuracy dưới $50\%$. Fine-tuning được kỳ vọng và cần được xác nhận bằng các bảng benchmark là chiến lược hiệu quả hơn do cho phép biểu diễn thích nghi với spectrogram RF. Trên dữ liệu thực tế, EfficientNet-B2, DINOv2 và SigLIP đạt tỷ lệ nhận diện trung bình lần lượt là $96,0\%$, $95,0\%$ và $90,5\%$. ConvNeXtV2 đạt $55,25\%$, còn ResNet50 đạt $18,5\%$.

Các kết quả chứng minh khả năng tổng quát hóa phụ thuộc vào từng backbone. EfficientNet-B2 đạt kết quả cao nhất và ổn định nhất giữa hai tập, vì vậy không có cơ sở kết luận Transformer luôn tốt hơn CNN. Sự khác biệt giữa benchmark và dữ liệu tự thu cần được phân tích cùng domain shift, khoảng cách thu và chất lượng spectrogram. Thí nghiệm tại 5~m, 10~m, 15~m và 20~m cho phép đánh giá mức duy trì khả năng nhận diện khi điều kiện truyền thay đổi. SNR thấp, clipping và stripe artifact có thể làm thay đổi đặc trưng đầu vào và ảnh hưởng trực tiếp đến dự đoán. Bên cạnh đó, độ dài cửa sổ, kích thước FFT, hop length, loại cửa sổ và phương pháp chuẩn hóa cần được khảo sát có kiểm soát vì chúng quyết định trực tiếp độ phân giải và hình học của ảnh đầu vào.

Những phát hiện của chương tạo cơ sở cho Chương~5 tổng kết đóng góp của đề tài, trình bày các hạn chế về dữ liệu và phạm vi đánh giá, đồng thời đề xuất hướng mở rộng hệ thống phát hiện drone trên nhiều điều kiện RF thực tế hơn.

\end{document}
